% SGA 3, Expose IV, English reconstruction.
% Scope: Lemma 1.7 through the end of section 2 (Remark 2.7), stopping
% immediately before the section 3 heading.
% Authority: Polo--Gille Exp4-13oct24.pdf, component-local pages 3--9
% (combined-reader pages 191--197).
% Jacob Reinhold's CC BY 4.0 Markdown lineage was used only as a phrasing
% comparison. The Polo--Gille PDF controls prose, formulas, numbering, notes,
% and diagrams.

% Source page: Exp4 component-local p. 3; combined-reader p. 191.
\medskip
\noindent\SGARefTarget{sga3:iv:lemma:1:7}\textbf{引理 1.7.}
\textit{沿用 \SGARefLink{sga3:iv:lemma:1:4}{1.4} 的记号,
若 \(u\) 和 \(v\) 是有效满态射, 且 \(v\) 是普遍满态射,
则 \(uv\) 是有效满态射.}

为此, 考虑下图
\[
\begin{array}{ccccc}
  S & \xleftarrow{u} & T & \leftleftarrows & T\times_S T \\
    &                 & \scriptstyle v\,\Big\uparrow
                      &                  & \scriptstyle v\times_Sv\,\Big\uparrow \\[-0.25em]
    &                 & U & \leftleftarrows & U\times_S U \\
    &                 & \substack{\uparrow\\[-0.6ex]\uparrow}
                      & \nearrow & \\
    &                 & U\times_T U & &
\end{array}
\tag{1.7.1}\label{sga3:iv:lemma:1:7:eq:1:7:1:01}
\]
由假设, 第一行与第一列正合; 此外, 根据
\SGARefLink{sga3:iv:corollary:1:5}{1.5} 和
\SGARefLink{sga3:iv:lemma:1:6}{1.6}, 由于 \(v\) 是普遍满态射,
\(v\times_Sv\) 是满态射. 由显然的图追踪即得结论.
事实上, 若 \(X(U)\) 的一个元素在 \(X(U\times_SU)\) 中的两个像相同,
则它在 \(X(U\times_TU)\) 中的两个像\textit{更是}相同,
从而由第一列的正合性, 它来自 \(X(T)\) 的一个元素.

% Source page: Exp4 component-local p. 4; combined-reader p. 192.
由于第一行正合, 只须验证上述元素在 \(X(T\times_ST)\) 中的两个像相同.
由于 \(v\times_Sv\) 是满态射, 只须验证这两个像在
\(X(U\times_SU)\) 中的像相同, 而这确实成立.

\medskip
\noindent\SGARefTarget{sga3:iv:proposition:1:8}\textbf{命题 1.8.}
\textit{考虑态射 \(U\xrightarrow{v}T\xrightarrow{u}S\).
若 \(u\) 和 \(v\) 是普遍有效满态射, 则 \(uv\) 是普遍有效满态射.
若 \(uv\) 是普遍有效满态射且 \(u\) 可平方,
则 \(u\) 是普遍有效满态射.}

第一个蕴含立即由 \SGARefLink{sga3:iv:lemma:1:7}{1.7} 得到.
对于第二个蕴含, 考虑 \SGARefLink{sga3:iv:proposition:1:8:eq:1:8:1:01}{下图}
(具有``双单纯''的形式):
\[
\begin{array}{ccccc}
  S & \xleftarrow{u} & T & \leftleftarrows & T\times_ST \\
  \scriptstyle uv\,\Big\uparrow
    & & \Big\uparrow & & \Big\uparrow \\[-0.25em]
  U & \longleftarrow & U\times_ST
    & \leftleftarrows & U\times_ST\times_ST \\
  \substack{\uparrow\\[-0.6ex]\uparrow}
    & & \substack{\uparrow\\[-0.6ex]\uparrow}
    & & \substack{\uparrow\\[-0.6ex]\uparrow} \\[-0.25em]
  U\times_SU & \longleftarrow & U\times_SU\times_ST
    & \leftleftarrows & U\times_SU\times_ST\times_ST .
\end{array}
\tag{1.8.1}\label{sga3:iv:proposition:1:8:eq:1:8:1:01}
\]
由 \(uv\) 为普遍有效满态射的假设, 第一、第二和第三列正合.
第二行正合, 因为 \(U\times_ST\to U\) 是有效满态射
(它在 \(U\) 上有截面); 同理, 第三行也正合.
由显然的图追踪可知第一行正合, 即 \(u\) 是有效满态射.
由于这些假设在任意基变换 \(S'\to S\) 下不变,
可知 \(u\) 实际上是普遍有效满态射.

\medskip
\noindent\SGARefTarget{sga3:iv:corollary:1:9}\textbf{推论 1.9.}
\textit{设 \(u:X\to Y\) 和 \(u':X'\to Y'\) 是普遍有效满态射,
并假设 \(Y\times Y'\) 存在. 则 \(X\times X'\) 存在, 且}
\[
  u\times u':X\times X'\longrightarrow Y\times Y'
\]
\textit{是普遍有效满态射.}

证明与 \SGARefLink{sga3:iv:corollary:1:5}{1.5} 相同, 使用下图
\[
\begin{array}{ccccc}
   & & Y' & \xleftarrow{u'} & X' \\
   & & \Big\uparrow & & \Big\uparrow \\[-0.25em]
  X & \longleftarrow & X\times Y' & \longleftarrow & X\times X' \\
  \scriptstyle u\,\Big\downarrow
    & & \Big\downarrow & \swarrow_{\scriptstyle u\times u'} & \\[-0.25em]
  Y & \longleftarrow & Y\times Y' & & .
\end{array}
\tag{1.9.1}\label{sga3:iv:corollary:1:9:eq:1:9:1:01}
\]

% Source page: Exp4 component-local p. 5; combined-reader p. 193.
\medskip
\noindent\SGARefTarget{sga3:iv:corollary:1:10}\textbf{推论 1.10.}
\textit{设 \(u:T\to S\) 是可平方态射, \(S'\to S\) 是一个
为普遍有效满态射的基变换态射. 则 \(u\) 为普遍有效满态射当且仅当}
\[
  u':T'=T\times_SS'\longrightarrow S'
\]
\textit{为普遍有效满态射:}
\[
\begin{array}{ccc}
  T & \xleftarrow{v'} & T' \\
  \scriptstyle u\,\Big\downarrow & & \Big\downarrow\,\scriptstyle u' \\[-0.25em]
  S & \xleftarrow{v} & S'.
\end{array}
\tag{1.10.1}\label{sga3:iv:corollary:1:10:eq:1:10:1:01}
\]
只须证明充分性. 若 \(u'\) 是普遍有效满态射, 则由
\SGARefLink{sga3:iv:proposition:1:8}{1.8}, \(vu'\) 也是普遍有效满态射.
由于 \(vu'=uv'\), 再次应用 \SGARefLink{sga3:iv:proposition:1:8}{1.8},
可知 \(u\) 是普遍有效满态射.

\medskip
\noindent\SGARefTarget{sga3:iv:remark:1:11}\textbf{注 1.11.}
\textit{同样的推理表明, 在 \SGARefLink{sga3:iv:corollary:1:10}{1.10} 中,
可将``普遍有效满态射''替换为``普遍满态射''、``既普遍又有效的满态射'',
或者仅仅替换为``满态射'' (在最后一种情形中, \(u\) 可平方的假设显然不必要).}

在 \SGARefLink{sga3:iv:proposition:1:8}{1.8} 的证明中,
我们使用了值得明确陈述的 \SGARefLink{sga3:iv:proposition:1:12}{下述结果}.

\medskip
\noindent\SGARefTarget{sga3:iv:proposition:1:12}\textbf{命题 1.12.}
\textit{设 \(u:T\to S\) 是有截面的态射. 则 \(u\) 是满态射;
若 \(T\times_ST\) 存在, 则它是有效满态射;
若此外 \(u\) 可平方, 则它是普遍有效满态射.}

第一个断言包含于 \SGARefLink{sga3:iv:lemma:1:4:item:a:01}{1.4(a)},
第三个断言立即由第二个得到, 因此只须证明第二个.
实际上, 还有更强的结论: 对每个不一定可表的函子
\(F:\mathscr C^{\circ}\to\mathbf{Set}\), 集合图
\[
  F(S)\longrightarrow F(T)\rightrightarrows F(T\times_ST)
\]
都是正合的. 这可视为 \v{C}ech 上同调形式体系的一个特例
(在维数 \(0\)!), 这里只作回顾. 仅假设 \(T\times_ST\) 存在, 并令
\[
  \check H^{0}(T/S,F)
  =\operatorname{Ker}\bigl(F(T)\rightrightarrows F(T\times_ST)\bigr).
\]

显然可将 \(\check H^{0}(T/S,F)\) 视为 \(\mathscr C_{/S}\) 中的变动对象
\(T\) 的反变函子: 每个 \(S\)-态射
\(T'\to T\) 都定义一个映射
\[
\tag{+}\label{sga3:iv:proposition:1:12:eq:unnamed:01}
  \check H^{0}(T/S,F)\longrightarrow\check H^{0}(T'/S,F).
\]
固定 \(\mathscr C_{/S}\) 中的 \(T,T'\). 一个熟知的计算表明,
若存在 \(S\)-态射 \(T'\to T\), 则相应的映射
\SGARefLink{sga3:iv:proposition:1:12:eq:unnamed:01}{(+)}
与该态射的选择无关.
\footnote{编者注: 下面是 M.~Demazure 告知我们的论证.
设 \(f,g:T'\to T\), \(\phi:T'\to T\times_ST\) 为分量是
\(f\) 和 \(g\) 的态射, 从而 \(p_1\circ\phi=f\) 及
\(p_2\circ\phi=g\). 于是 \(F(\phi):F(T\times_ST)\to F(T')\) 满足
\(F(\phi)\circ F(p_1)=F(f)\) 及 \(F(\phi)\circ F(p_2)=F(g)\).
然而, 对每个 \(x\in\check H^0(T/S,F)\), 有
\(F(p_1)(x)=F(p_2)(x)\). 对等式两边应用 \(F(\phi)\), 得到
\(F(f)(x)=F(g)(x)\), 这说明 \(f\) 和 \(g\) 诱导相同的映射.}

% Source page: Exp4 component-local p. 6; combined-reader p. 194.
因此, \(\check H^{0}(T/S,F)\) 可视为与集合
\(\operatorname{Ob}(\mathscr C_{/S})\) 关联的范畴上的函子,
其中以支配关系赋予该集合预序: 若存在 \(S\)-态射 \(T'\to T\),
则称 \(T'\) 支配 \(T\). 特别地, 若 \(T\) 和 \(T'\) 在后一个范畴中
同构, 即两者互相支配, 则
\SGARefLink{sga3:iv:proposition:1:12:eq:unnamed:01}{(+)} 是集合的同构.
特别地, 这适用于 \(T'\) 是 \(\mathscr C_{/S}\) 的终对象的情形,
其实质上就是 \(S\) 自身. 对象 \(T\) 总是支配 \(T'=S\),
反之成立恰好当 \(T/S\) 有截面. 这就证明了
\SGARefLink{sga3:iv:proposition:1:12}{1.12} 所宣称的较强形式.

\medskip
\noindent\SGARefTarget{sga3:iv:remark:1:13}\textbf{注 1.13.}
\textit{对于各种应用, 本篇引入的概念和陈述的结果需要进一步推广到
有相同目标的态射族 \(u_i:T_i\to S\), 而不限于单个态射
\(u:T\to S\). 若对 \(\mathscr C\) 的每个对象 \(X\), 相应的映射}
\[
  X(S)\longrightarrow\prod_i X(T_i)
\]
\textit{都是单射, 则称这样的态射族为\emph{满的 (epimorphic)}.
类似地可引入有效满态射族的概念, 以及这些概念的普遍形式.
在下文需要时, 将不再说明地使用本篇结果在这一更一般情形下的推广.}

\medskip
\noindent\SGARefTarget{sga3:iv:remark:1:14}\textbf{注 1.14.}
\textit{既是单态射又是有效满态射的每个态射都是同构.}

事实上, 沿用 \SGARefLink{sga3:iv:definition:1:3}{1.3} 的记号,
\(T\to S\) 是单态射意味着两个态射
\[
  \operatorname{pr}_1,\operatorname{pr}_2:T\times_ST
  \rightrightarrows T
\]
相等 (并且都是同构). 而图
\[
  C\overset{v}{\underset{v}{\rightrightarrows}}B\xrightarrow{u}A
\]
正合仅当 \(u\) 为同构, 这立即由
\SGARefLink{sga3:iv:definition:1:3}{定义} 得到.
\footnote{编者注: 既是单态射又是满态射的态射未必是同构.
例如, 在 \(\mathscr C=\mathbf{Sch}\) 中, 态射
\[
  \operatorname{Spec}(\mathbf F_p)\amalg
  \operatorname{Spec}(\mathbf Z[1/p])
  \longrightarrow\operatorname{Spec}(\mathbf Z)
\]
是单态射, 也是满射的满态射, 却不是同构.}

\SGARefTarget{sga3:iv:section:2}\section*{2. 下降态射}
\addcontentsline{toc}{section}{2. 下降态射}

回顾以下定义.

\medskip
\noindent\SGARefTarget{sga3:iv:definition:2:1}\textbf{定义 2.1.}
\textit{设 \(f:S'\to S\) 是使 \(S''=S'\times_SS'\) 存在的态射,
\(u':X'\to S'\) 是 \(S'\) 上的一个对象.
\(X'/S'\) 上相对于 \(f\) 的一个\emph{粘合数据 (gluing datum)},
是指一个 \(S''\)-同构}
\[
  c:X''_1\xrightarrow{\sim}X''_2,
\]
\textit{其中, 对 \(i=1,2\), 用 \(X''_i\) 表示 \(X'/S'\)
在投影 \(\operatorname{pr}_i:S''\to S'\) 下的逆像, 并假设该逆像存在.
若粘合数据 \(c\) 满足``余循环条件''}
\[
  \operatorname{pr}_{3,1}^{*}(c)
  =\operatorname{pr}_{3,2}^{*}(c)\,
   \operatorname{pr}_{2,1}^{*}(c),
\]
\textit{则称之为一个\emph{下降数据 (descent datum)}.
其中, 对 \(1\leq j<i\leq3\), \(\operatorname{pr}_{i,j}\)
是从}
\[
  S'''=S'\times_SS'\times_SS'
\]
\textit{到 \(S''\) 的典范投影 (此时也假设 \(S'''\) 存在);
\(\operatorname{pr}_{i,j}^{*}(c)\) 是 \(c\) 的逆像,
视为从 \(X'''_j\) 到 \(X'''_i\) 的 \(S'''\)-态射;
对从 \(1\) 到 \(3\) 的每个整数 \(k\), 用 \(X'''_k\) 表示
\(X'/S'\) 在第 \(k\) 个投影 \(q_k:S'''\to S'\) 下的逆像,
并假设该逆像存在.}

% Source page: Exp4 component-local p. 7; combined-reader p. 195.
在 \SGARefLink{sga3:iv:definition:2:1}{定义} 的第二部分中,
我们因而使用了通常的等同和记号滥用,
\footnote{编者注: 例如, 我们一方面作了等同
\(\operatorname{pr}_{2,1}^{*}(X''_1)
 =\operatorname{pr}_{3,1}^{*}(X''_1)\), 另一方面作了等同
\(\operatorname{pr}_{2,1}^{*}(X''_2)
 =\operatorname{pr}_{3,2}^{*}(X''_2)\).}
经验表明这些做法没有问题, 但在严格阐述下降理论时显然必须避免;
下降理论本身在某种程度上也必须为这些习惯性的术语滥用提供依据.
Giraud 已写出这样的形式化阐述 (\SGARefLink{sga3:iv:bib:d}{[D]}),
目的是为 \SGARefLink{sga3:iv:bib:sga:1}{SGA~1} 中从未写成的
第 VII 篇提供依据并使之精确. 关于本讨论班反复使用的忠实平坦下降结果的
详细说明, 可参阅 \SGARefLink{sga3:iv:bib:sga:1}{SGA~1} 第 VIII 篇.

仍设 \(f:S'\to S\) 为使 \(S''=S'\times_SS'\) 存在的态射,
\(X\) 为 \(S\) 上的对象, 满足 \(X'=X\times_SS'\) 和
\(X''=X\times_SS''\) 存在. 于是 \(X'\) 在
\(\operatorname{pr}_i\), \(i=1,2\), 下的逆像存在且典范同构;
因此, \(X'/S'\) 带有相对于 \(f\) 的典范粘合数据.
当 \(S'''\) 和 \(X'''=X\times_SS'''\) 存在时,
这实际上是一个下降数据. 若 \(Y\) 是 \(S\) 上另一个满足与
\(X/S\) 相同条件的对象, 则对每个 \(S\)-态射 \(X\to Y\),
相应的 \(S'\)-态射 \(X'\to Y'\) 与 \(X'\) 和 \(Y'\) 上的
典范粘合数据相容. 特别地, 若 \(S'\to S\) 可平方, 则
\[
  X\longmapsto X'=X\times_SS'
\]
是从 \(\mathscr C_{/S}\) 到``\(S'\) 上带有相对于 \(f\) 的下降数据的
对象范畴''的函子; 这个范畴的定义留给读者,
它是``\(S'\) 上带有相对于 \(f\) 的粘合数据的对象范畴''的全子范畴.

在此约定下, 给出 \SGARefLink{sga3:iv:definition:2:2}{如下定义}.

\medskip
\noindent\SGARefTarget{sga3:iv:definition:2:2}\textbf{定义 2.2.}
\textit{态射 \(f:S'\to S\) 称为\emph{下降态射}
(分别地, \emph{有效下降态射}), 是指 \(f\) 可平方
(即对每个 \(X\to S\), \(X\times_SS'\) 都存在),
且上述函子}
\[
  X\longmapsto X'=X\times_SS'
\]
\textit{从 \(S\) 上的对象范畴 \(\mathscr C_{/S}\) 到
\(S'\) 上带有相对于 \(f\) 的下降数据的对象范畴,
是全忠实的 (分别地, 是范畴等价).}

第一个概念仅用粘合数据即可表述, 因而无须使用三重纤维积 \(S'''\):
\(f\) 是下降态射恰好当它可平方, 且 \(X\mapsto X'\)
是从 \(\mathscr C_{/S}\) 到 \(S'\) 上带有相对于 \(f\) 的粘合数据的
对象范畴的全忠实函子. 展开来说, 这意味着对 \(S\) 上的任意两个对象
\(X,Y\), 下面的集合 \SGARefLink{sga3:iv:definition:2:2:eq:x:01}{图} 正合:
\[
\tag{x}\label{sga3:iv:definition:2:2:eq:x:01}
  \operatorname{Hom}_S(X,Y)
  \longrightarrow \operatorname{Hom}_{S'}(X',Y')
  \overset{p_1}{\underset{p_2}{\rightrightarrows}}
  \operatorname{Hom}_{S''}(X'',Y'').
\]

% Source page: Exp4 component-local p. 8; combined-reader p. 196.
这里, 对 \(i=1,2\), \(p_i(u')\) 表示
\(u'\in\operatorname{Hom}_{S'}(X',Y')\) 在投影
\(\operatorname{pr}_i:S''\to S'\) 下的逆像. 事实上,
按照定义, 箭头对 \((p_1,p_2)\) 的核恰好是与典范粘合数据相容的
\(S'\)-态射 \(X'\to Y'\) 组成的集合.

由逆像 \(Y'\) 和 \(Y''\) 的定义, 有以下典范双射
\[
  \operatorname{Hom}_{S'}(X',Y')\simeq\operatorname{Hom}_S(X',Y),
  \qquad
  \operatorname{Hom}_{S''}(X'',Y'')\simeq\operatorname{Hom}_S(X'',Y).
\]
因此 \SGARefLink{sga3:iv:definition:2:2:eq:x:01}{(x)} 正合等价于
\[
\tag{xx}\label{sga3:iv:definition:2:2:eq:xx:01}
  \operatorname{Hom}_S(X,Y)
  \longrightarrow\operatorname{Hom}_S(X',Y)
  \rightrightarrows\operatorname{Hom}_S(X'',Y),
\]
正合; 后者是将 \(\operatorname{Hom}_S(-,Y)\) 应用于
\(\mathscr C_{/S}\) 中的图
\[
\tag{xxx}\label{sga3:iv:definition:2:2:eq:xxx:01}
  X''\rightrightarrows X'\longrightarrow X,
\]
所得到的, 而这个图本身由
\[
  S''\rightrightarrows S'\longrightarrow S
\]
经基变换 \(X\to S\) 得到. 根据 \SGARefLink{sga3:iv:lemma:1:6}{1.6},
这证明了 \SGARefLink{sga3:iv:proposition:2:3}{下面命题} 的第一部分.

\medskip
\noindent\SGARefTarget{sga3:iv:proposition:2:3}\textbf{命题 2.3.}
\textit{设 \(f:S'\to S\) 是态射. 若 \(f\) 是普遍有效满态射,
则它是下降态射. 反之, 若 \(S\) 是 \(\mathscr C\) 的可平方对象,
即它与 \(\mathscr C\) 的每个对象的乘积都存在,
则每个下降态射都是普遍有效满态射.}

只需证明逆命题. 若 \(f\) 是下降态射, 须证明它是普遍有效满态射:
对每个态射 \(X\to S\), 图
\SGARefLink{sga3:iv:definition:2:2:eq:xxx:01}{(xxx)} 必须正合,
或等价地, 对 \(\mathscr C\) 的每个对象 \(Z\),
它在 \(\operatorname{Hom}(-,Z)\) 下的像必须是正合的集合图.
由假设, \(Z\times S\) 存在. 设 \(Y\) 是它定义的
\(\mathscr C_{/S}\) 的对象. 于是
\SGARefLink{sga3:iv:definition:2:2:eq:xxx:01}{(xxx)} 在
\(\operatorname{Hom}(-,Z)\) 下的像同构于其在
\(\operatorname{Hom}_S(-,Y)\) 下的像, 而后者由关于 \(f\) 的假设正合.

因此, 关于下降态射的结果可以应用于普遍有效满态射. 例如:

\medskip
\noindent\SGARefTarget{sga3:iv:proposition:2:4}\textbf{命题 2.4.}
\textit{设 \(f:S'\to S\) 是下降态射 (例如普遍有效满态射). 则:}
\begin{enumerate}
  \item[(a)]\SGARefTarget{sga3:iv:proposition:2:4:item:a:01}
    \textit{对每个 \(S\)-态射 \(u:X\to Y\), 态射 \(u\) 为同构
    (分别地, 单态射) 当且仅当 \(u':X'\to Y'\) 为同构 (分别地, 单态射).}
  \item[(b)]\SGARefTarget{sga3:iv:proposition:2:4:item:b:01}
    \textit{设 \(X,Y\) 是 \(S\) 的两个子对象, \(X',Y'\) 是它们在
    \(S'\) 中的逆像. 则 \(X\) 包含于 \(Y\) (分别地, 等于 \(Y\))
    当且仅当 \(X'\) 和 \(Y'\) 具有相同的关系.}
\end{enumerate}

% Source page: Exp4 component-local p. 9; combined-reader p. 197.
对于 \SGARefLink{sga3:iv:proposition:2:4:item:a:01}{(a)},
由 \SGARefLink{sga3:iv:definition:2:2}{定义} 可知,
若 \(u'\) 在带粘合数据的对象范畴中是同构, 则 \(u\) 是同构.
而 \(S'\) 上与粘合数据相容的对象同构, 总是带粘合数据的对象的同构:
其逆也与粘合数据相容. 对于
\SGARefLink{sga3:iv:proposition:2:4:item:b:01}{(b)},
只须证明: 若 \(X'\) 包含于 \(Y'\), 即存在 \(S'\)-态射
\(X'\to Y'\), 则 \(S\) 上的 \(X,Y\) 也具有同样性质.
由于 \(Y'\to S'\), 因而也包括 \(Y''\to S''\), 是单态射,
态射 \(X'\to Y'\) 自动与粘合数据相容, 从而来自一个
\(S\)-态射 \(X\to Y\).

这个证明更一般地适用于 \(S\) 上的两个对象 \(X,Y\),
其中 \(Y\to S\) 是单态射, 而要判断 \(X\to S\) 是否经过 \(Y\) 分解:
只须 \(X'\to S'\) 经过 \(Y'\) 分解即可.

\medskip
\noindent\SGARefTarget{sga3:iv:corollary:2:5}\textbf{推论 2.5.}
\textit{设 \(f:S'\to S\) 是普遍有效满态射, \(g:S\to T\)
是使 \(S\times_TS\) 存在的态射. 假设 \(S''=S'\times_SS'\)
也是 \(S'\) 在 \(T\) 上与自身的纤维积, 即典范态射}
\[
  S'\times_SS'\xrightarrow{\sim}S'\times_TS'
\]
\textit{为同构. 则 \(g:S\to T\) 是单态射 (反之当然也成立).}

事实上, 考虑 笛卡尔 图
\[
\begin{array}{ccc}
  S'\times_SS' & \xrightarrow{\sim} & S'\times_TS' \\
  \Big\downarrow & & \Big\downarrow\,\scriptstyle f \\[-0.25em]
  S & \longrightarrow & S\times_TS ,
\end{array}
\tag{2.5.1}\label{sga3:iv:corollary:2:5:eq:2:5:1:01}
\]
其中第二个水平箭头是对角态射. 由
\SGARefLink{sga3:iv:corollary:1:9}{1.9}, 第二个竖直箭头 \(f\)
是普遍有效满态射; 由假设, 第一个水平箭头是同构.
由 \SGARefLink{sga3:iv:proposition:2:4:item:a:01}{2.4(a)},
或者同样由 \SGARefLink{sga3:iv:proposition:2:4:item:b:01}{2.4(b)},
\footnote{编者注: 应用于 \(f:S'\times_TS'\to S\times_TS=Y\).}
可知 \(S\to S\times_TS\) 也是同构, 这正表示
\(g:S\to T\) 是单态射.

\medskip
\noindent\SGARefTarget{sga3:iv:remark:2:6}\textbf{注 2.6.}
\textit{在 \SGARefLink{sga3:iv:definition:2:1}{2.1} 中用某些逆极限之间
的态射引入的概念, 显然可用所考虑的对象 \(S,S',X'\) 定义的反变函子表述.
在 \SGARefLink{sga3:iv:definition:2:1}{2.1} 所考虑的纤维积存在的条件下,
\(X'/S'\) 上相对于 \(S'\to S\) 的粘合数据 (分别地, 下降数据),
与下列范畴中相应对象的粘合数据 (分别地, 下降数据) 一一对应:}
\[
  \widehat{\mathscr C}
  =\operatorname{Hom}(\mathscr C^{\circ},\mathbf{Set}).
\]
\textit{因此, 如有需要, 可将这些概念推广到不对
\(\mathscr C\) 中的纤维积施加存在性假设的情形.}

\medskip
\noindent\SGARefTarget{sga3:iv:remark:2:7}\textbf{注 2.7.}
\textit{\SGARefLink{sga3:iv:section:2}{本节}引入的概念可推广到态射族.
借助筛的概念 (\SGARefLink{sga3:iv:subsection:4:1}{4.1}),
也可将它们表述得更内在. 关于这些问题, 请读者参阅
\SGARefLink{sga3:iv:bib:d}{[D]}.}
% Translation statistics: source_lines=464; output_lines=444; inline_math=242; display_math=26; footnotes=4; pending_TODO=0.
% Review: complete assigned component; mathematical symbols, targets, references, tags, environments and footnotes checked; formula prose translated.
